\chapter{LITERATURE REVIEW}\label{ch:lit}

\begin{doublespacing}

Multi-robot search and rescue draws on three largely independent
research traditions: online competitive search theory, market-based
multi-robot task allocation, and energy-constrained graph
exploration. This chapter reviews each in turn, identifies the gap
they leave when taken together, and positions the present work
inside that gap. A reader without a background in any of the three
may find it useful to read the \emph{primer} at the start of each
section for a plain-language explanation before moving on to the
formal review.

\section{Competitive Search Theory}

\paragraph{Primer.} The field of \emph{online algorithms} studies
decisions that have to be made with incomplete information. A
classic textbook example is the following. You are standing at the
origin of an infinite straight road. A cow you are looking for is
somewhere on this road, either to your left or to your right, at an
unknown distance. You can walk in either direction and will
recognize the cow when you reach it, but you cannot see her from
afar. How should you search so that you do not walk much farther
than strictly necessary? If you always happen to guess the correct
direction first, you walk exactly the distance to the cow; if you
guess wrong, you have to backtrack. ``Competitive search theory''
proves how good any strategy can possibly be in the worst case,
measured as a ratio to the distance an all-knowing agent would
walk. This ratio is called the \emph{competitive ratio}. It is the
same idea the Finder Competitive Ratio (FCR) in this thesis
captures, but in two dimensions and with many searchers.

Competitive analysis measures an online algorithm by the worst-case
ratio of its cost to the cost an omniscient offline algorithm would
incur on the same instance~\citep{Sleator1985}. The canonical
instantiation of this idea for search is the \emph{cow-path problem}:
a searcher stands at the origin of an infinite half-line, a target is
somewhere on the line, and the cost is the distance traveled before
the target is found. \citet{BaezaYates1993} proved that a geometric doubling strategy
achieves a deterministic competitive ratio of 9 and that no
deterministic strategy can do better. \citet{Kao1996} showed that
randomization brings the ratio down to approximately $4.59$. These results are important because they
characterize the intrinsic cost of not knowing the target's location,
independent of the specific algorithm chosen.

Extensions to parallel search introduce multiple agents and ask how
the ratio changes. \citet{LopezOrtiz2004} gave competitive ratios for
$p$ parallel searchers exploring $m$ rays emanating from a common
origin. \citet{Chrobak2015} studied group search on a line and proved a
striking dichotomy: with wireless communication, two agents achieve
ratio 3, while with only face-to-face (rendezvous) communication the
ratio remains the single-searcher 9. That is, the \emph{mode} of
communication matters more than the number of searchers. This
observation motivates the first finding of the present work: that
coordination alone, not the number of robots, accounts for the bulk of
the improvement over random search. The competitive-ratio literature
has, however, focused on one-dimensional geometries and has not
integrated energy constraints.

A rigorous theoretical foundation for search in more general settings
is provided by \citet{AlpernGal2003}, whose book \emph{The Theory of
Search Games and Rendezvous} unifies the competitive-ratio and
game-theoretic approaches for search on graphs, lines, and compact
spaces. \citeauthor{AlpernGal2003} formalize the minimax and expected-cost
objectives for search problems and establish that competitive ratio
is the natural criterion when the adversary may place the target
adversarially. The present thesis adopts this perspective but extends
it to the probabilistic setting (where the target is drawn from a known
prior rather than placed adversarially) and to the two-dimensional,
energy-constrained, multi-robot regime that
\citeauthor{AlpernGal2003} did not consider.

\section{Market-Based Multi-Robot Coordination}

\paragraph{Primer.} How should several robots divide a pile of
tasks? One approach is to have a central computer solve an
optimization problem and hand each robot its assignment. A different
approach, inspired by how markets work in economics, is to let the
robots \emph{bid} for tasks. Each robot computes how much it
``wants'' each task (based on how close the task is, how hard it
looks, how well the robot is suited for it, and so on), the robots
exchange bids with each other, and the highest bidder wins each
task. Market-based coordination is popular in robotics because it
is simple, it scales to large teams without a bottleneck, and it
gracefully tolerates robots dropping in and out of the team. The
particular auction mechanism this thesis uses is \emph{sequential
single-item (SSI)} auctions: tasks are sold one at a time, the
highest bidder wins each one, and the process repeats until all
robots have assignments for the current round. SSI auctions are a
well-studied compromise between the simplicity of a one-shot
assignment and the flexibility of a fully iterative market.

\citet{Gerkey2004} introduced the multi-robot
task allocation (MRTA) taxonomy that classifies problems by three
axes: single-task versus multi-task robots, single-robot versus
multi-robot tasks, and instantaneous versus time-extended assignment.
In this taxonomy, the present problem is
\emph{single-task, multi-robot, instantaneous} (ST-MR-IA): each site
is a single task, multiple robots may claim different sites in the
same round, and the assignment is recomputed every round.

Auction-based task allocation was popularized by the TraderBots
architecture of \citet{Dias2006} and the earlier market economy of
\citet{Zlot2002}. \citet{Koenig2006}
analyzed sequential single-item (SSI) auctions---in which items are
auctioned one at a time and the globally highest bidder wins each
item---and proved constant-factor approximation guarantees against
centralized optima for several objective functions. SSI auctions are
attractive because they are simple, communication-efficient, and
produce provably near-optimal assignments. This thesis adopts SSI as
the basic coordination primitive.

A comprehensive systematic review by \citet{ACMSurveyMRTA2025}
examined 78 MRTA papers published between 2013 and 2024 and confirmed
that market-based methods remain the dominant paradigm in multi-robot
coordination, outperforming both centralized optimization and
learning-based approaches on dynamic, time-extended tasks. The review
identifies energy-aware task allocation as one of the principal open
problems: the overwhelming majority of surveyed papers assume unlimited
robot mobility, and none incorporate per-sortie energy budgets coupled
with probabilistic site priors. This gap directly motivates the present
work.

A critical practical dimension is what happens when the communication
channel degrades. \citet{OtteKuhlman2020} compared six auction variants
experimentally as communication quality ranges from perfect to
nonexistent, finding that different auction protocols degrade at very
different rates and that the choice of auction matters far more than the
bid function in degraded-radio settings. Their work provides the
theoretical footing for an important limitation of the present thesis:
Models M2, M3, and M4 all assume reliable at-base communication,
and their performance guarantees do not extend to dropped-radio
scenarios. The sensitivity of M4*'s FCR advantage to communication
degradation is a direct open problem identified in Chapter~VI.

The market-based literature has, however, historically studied mobility
as unconstrained: a robot is expected to be able to reach any assigned
task, which does not match the energy-constrained sortie regime that
this thesis studies.

\section{Energy-Constrained and Probabilistic Search}

\paragraph{Primer.} The two ideas combined in this section are, on
the one hand, that a real robot runs out of battery, and on the
other, that not every location is equally likely to hold what the
robot is searching for. The energy-constrained branch of robotics
research treats the battery as a hard constraint: a robot must
plan every excursion knowing that at some point it has to return
to base to recharge, which fundamentally changes what strategies
are available. The probabilistic branch treats the target
location as a random variable---some sites are more likely to
contain it than others---and asks how the robot should allocate its
search effort across the possibilities. A famous classical result
in this area is \emph{Stone's rule}~\citep{Stone1975}: optimal
search effort should be allocated in proportion to (probability
of target being there) divided by (cost of searching there). This
is the direct ancestor of the $p/d$ bid function in this thesis,
with $p$ the probability and $d$ the search cost (distance). The
novel $p/d^2$ bid that the thesis recommends is a modification of
Stone's rule tailored to the situation where battery is the binding
constraint.

The third strand of related work concerns search under finite energy
and under probabilistic priors. \citet{Betke1995} introduced \emph{piecemeal learning}, in which a
robot must periodically return to the origin and therefore can only
execute bounded-length excursions. \citet{Das2015} studied
energy-constrained depth-first search
and proved that the number of robots required to explore a graph
grows with the ratio of graph diameter to battery capacity. The common
thread is that energy constraints fundamentally alter the structure of
optimal search strategies; a fuel-limited robot cannot use the same
plans as a fuel-unlimited one.

On the probabilistic side, Stone's classical \emph{Theory of Optimal
Search}~\citep{Stone1975} established that search effort should be
allocated in proportion to the ratio of target probability to search
cost. This is the theoretical basis for the $p_i / d$ bid function
used in this thesis: it is Stone's rule evaluated locally at each
decision point. \citet{Bourgault2002} applied information-theoretic
criteria (specifically mutual information) to coordinated target
search and demonstrated experimentally that information-greedy
policies outperform distance-greedy policies. The observation
generalizes to the present work.

An important contemporary alternative to analytical bid functions is
deep reinforcement learning (DRL). \citet{KimKim2025} used a
market-based strategy combined with Fisher information maximization and
weighted Voronoi partitioning to coordinate multi-robot search in open
fields, achieving competitive performance on continuous probability
maps. The Fisher-information criterion is richer than the $p/d$ or
$p/d^2$ heuristic in that it accounts for full posterior covariance,
but it comes at substantially higher computational cost and does not
admit closed-form competitive-ratio bounds. The present thesis occupies
a complementary point in the design space: closed-form bid functions
and analytical bounds are traded for the richer posterior updates of
information-theoretic methods, but what is gained is exact formal
guarantees and a one-line deployment prescription.

A structural robustness result directly relevant to this thesis is
provided by \citet{AutomaticaAuctions2023}, who derived intervals on
each cost parameter within which perturbations leave the SSI auction's
output allocation unchanged. This sensitivity analysis provides the
theoretical machinery to formally quantify how stable the M4* bid
function is to perturbations in site probabilities---a natural
extension identified in Chapter~VI as future work.

\section{Gap and Positioning}

Taken together, the three literatures leave a gap. Competitive-ratio
analysis gives worst-case bounds but assumes unrealistic
one-dimensional geometries and does not model energy. Market-based
coordination scales to many robots but assumes unbounded mobility.
Energy-constrained exploration models batteries but does not use
auctions or compute competitive ratios. No prior work unifies all
three into a single framework that asks: what is the expected
competitive ratio of an auction-coordinated multi-robot search team
under a finite energy budget?

The present thesis fills that gap. Four novel contributions emerge
from it that do not appear in the surveyed literature.

First, the \emph{Finder Competitive Ratio} (FCR) metric is
introduced as a formal performance measure for the target-finding
problem. Existing competitive-ratio analyses measure the worst-case
team cost; FCR instead measures the travel of the specific robot that
discovers the target, normalized by the omniscient-optimum distance.
This finder-centric perspective captures operationally relevant
performance---the time from deployment to discovery---and is not
equivalent to minimizing total-team distance, as the Target-Delay
Anomaly of Chapter~V demonstrates.

Second, a sortie model is formalized in which each energy budget $E$
constrains a single round trip, not a whole mission. This contrasts
with prior energy models that either restrict total lifetime
travel~\citep{Betke1995} or graph traversal under a global
budget~\citep{Das2015}. The sortie model matches the battery-swap
reality of UAV deployments and enables the multi-site chaining
strategy of Model~4.

Third, the thesis derives a \emph{cost-foreclosure argument} that
justifies a superlinear distance penalty in the bid function. No prior
work in auction-based coordination has analyzed how the choice of bid
exponent interacts with a per-sortie energy cap.

Fourth, a $p_i / d^2$ bid function is identified empirically as
optimal under the sortie model, beating the Stone-rule-derived $p_i/d$
bid. The advantage persists under prior misspecification noise and
across three non-Euclidean cost variants, demonstrating robustness
that the analytical argument predicts but does not prove.

In addition, the thesis conducts an ablation that disentangles two
effects previously conflated in the MRTA literature: the value of
probability-aware bids versus the value of the sequential iterative
allocation mechanism.
The ablation reveals that a probability-weighted Hungarian assignment
H$_{p/d}$ performs \emph{worse} than a distance-only Hungarian baseline
H$_d$, while the sequential single-item auction M3 outperforms both.
This establishes that the iterative site-removal structure of the SSI
auction---not probability-aware bid values alone---is the essential
contributor to coordination quality under reliable communication.
The result concretely answers the question implicitly raised by
\citet{OtteKuhlman2020} of whether coordination structure or bid
design is the more impactful design decision.

\end{doublespacing}
